\usepackage[T1]{fontenc}
\usepackage[utf8]{inputenc}
\usepackage{lmodern}
\usepackage[a4paper,inner=26mm,outer=24mm,top=25mm,bottom=26mm,headheight=14pt]{geometry}
\usepackage{amsmath,amssymb,amsthm,mathtools,bm}
\usepackage{microtype}
\usepackage{booktabs,longtable,array,tabularx}
\usepackage{graphicx}
\usepackage[dvipsnames]{xcolor}
\usepackage{tikz}
\usetikzlibrary{arrows.meta,positioning,calc,decorations.pathreplacing}
\usepackage{xurl}
\usepackage{fancyhdr}
\usepackage{makeidx}
\makeindex
\usepackage{hyperref}
\usepackage{bookmark}
\hypersetup{colorlinks=true,linkcolor=MidnightBlue,urlcolor=MidnightBlue,citecolor=MidnightBlue,
 pdftitle={Haagerup Symmetry and Conformal Field Theory},
 pdfsubject={A graduate textbook on fusion categories, the central-charge-eight construction, and the Haagerup chain}}
\urlstyle{tt}
\setlength{\emergencystretch}{2em}
\setlength{\parskip}{0.2em}
\pagestyle{fancy}
\fancyhf{}
\fancyhead[L]{\small\nouppercase{\leftmark}}
\fancyfoot[C]{\thepage}
\renewcommand{\headrulewidth}{0.3pt}
\makeatletter
\renewcommand{\chaptermark}[1]{\markboth{\if@mainmatter\thechapter.\ \fi#1}{}}
\makeatother
\newtheorem{theorem}{Theorem}[chapter]
\newtheorem{lemma}[theorem]{Lemma}
\newtheorem{proposition}[theorem]{Proposition}
\newtheorem{corollary}[theorem]{Corollary}
\theoremstyle{definition}
\newtheorem{definition}[theorem]{Definition}
\newtheorem{example}[theorem]{Worked example}
\newtheorem{exercise}{Exercise}[chapter]
\newtheorem{constructioninput}[theorem]{Construction input}
\newtheorem{numericalrecord}[theorem]{Numerical record}
\let\textbookexample\example
\let\endtextbookexample\endexample
\RenewDocumentEnvironment{example}{o}{%
 \IfNoValueTF{#1}{\textbookexample}{\textbookexample[#1]}%
 \addcontentsline{loe}{workedexample}{%
  \protect\numberline{\thetheorem}%
  \IfNoValueTF{#1}{Worked example}{#1}}%
}{\endtextbookexample}
\makeatletter
\newcommand{\l@workedexample}{\@dottedtocline{1}{0em}{3.8em}}
\newcommand{\toclevel@workedexample}{1}
\newcommand{\listofworkedexamples}{%
 \chapter*{Worked examples}%
 \phantomsection
 \addcontentsline{toc}{chapter}{Worked examples}%
 \markboth{Worked examples}{Worked examples}%
 \@starttoc{loe}}
\makeatother
\theoremstyle{remark}
\newtheorem{remark}[theorem]{Remark}
\newenvironment{solution}{\begin{proof}[Solution]}{\end{proof}}
\newcommand{\C}{\mathbb C}
\newcommand{\R}{\mathbb R}
\newcommand{\Q}{\mathbb Q}
\newcommand{\Z}{\mathbb Z}
\newcommand{\N}{\mathbb N}
\newcommand{\one}{\mathbf1}
\newcommand{\HH}{\mathcal H}
\providecommand{\mathscr}[1]{\mathcal{#1}}
\newcommand{\CC}{\mathbb C}
\newcommand{\QQ}{\mathbb Q}
\newcommand{\ZZ}{\mathbb Z}
\DeclareMathOperator{\Rep}{Rep}
\DeclareMathOperator{\Com}{Com}
\DeclareMathOperator{\FPdim}{FPdim}
\DeclareMathOperator{\qdim}{qdim}
\DeclareMathOperator{\wt}{wt}
\DeclareMathOperator{\Tr}{Tr}
\DeclareMathOperator{\End}{End}
\DeclareMathOperator{\Hom}{Hom}
\DeclareMathOperator{\Aut}{Aut}
\DeclareMathOperator{\Span}{span}
\DeclareMathOperator{\Vir}{Vir}
\DeclareMathOperator{\Irr}{Irr}
\DeclareMathOperator{\ch}{ch}
\DeclareMathOperator{\rank}{rank}
\DeclareMathOperator{\diag}{diag}
\DeclareMathOperator{\Spec}{Spec}
\renewcommand{\Vec}{\operatorname{Vec}}
\newcommand{\sourcefile}[1]{\texttt{\nolinkurl{#1}}}
\newcommand{\norm}[1]{\left\lVert#1\right\rVert}
\newcommand{\inner}[2]{\left\langle#1,#2\right\rangle}
\newcommand{\chapterinput}[1]{\input{chapters/#1.tex}}
